\documentclass{article}
\usepackage[left=2cm, right=2cm, top=2cm, bottom=2cm]{geometry}
\usepackage{graphicx}
\usepackage{amsmath}
\usepackage{amssymb}
\usepackage{amsthm}
\usepackage{fancyhdr}
\usepackage{verbatim}
\usepackage{listings}
\usepackage{xcolor}
\usepackage{pdfpages}
\usepackage{cancel}
\usepackage{physics}
\usepackage{esint}
\usepackage{braket}

\lstset{
    language=C++,
    basicstyle=\ttfamily\footnotesize,
    keywordstyle=\color{blue},
    commentstyle=\color{green},
    stringstyle=\color{red},
    numbers=left,
    numberstyle=\tiny\color{gray},
    frame=single,
    breaklines=true,
    captionpos=b,
}


\title{MTH 446: Lecture \# 4}
\author{Cliff Sun}

\newtheorem{theorem}{Theorem}[section]
\newtheorem{lemma}[theorem]{Lemma}
\newtheorem{definition}[theorem]{Definition}
\newtheorem{conjecture}[theorem]{Conjecture}
\newtheorem{proposition}[theorem]{Proposition}
\newtheorem{corollary}[theorem]{Corollary}
\newtheorem{one minute paper}[theorem]{One Minute Paper}

\pagestyle{fancy}
\lhead{\textbf{\thepage}\ \ \nouppercase{\rightmark}}
\chead{MTH 446: Lecture \# 4}
\rhead{Cliff Sun}

\begin{document}

\maketitle

\section*{Lecture Span}
\begin{enumerate}
    \item Roots of complex numbers
\end{enumerate}

\section*{Simple multiplication example}

In general, evaluating if $z_1z_2 = z_3$, we can either (1) distribute and evalaute element-by-element or (2) check that the modulus of the RHS = LHS and that the arguments of both are the same.

\textbf{Need to review what values of cos correspond to $\theta$ values}. 

To obtain the argument of a complex number, calculate 

\begin{equation}
    z = |z|\left( \frac{z}{|z|} \right)
\end{equation}

Here, then the complex number $z/|z|$ corresponds to some $e^{i\theta}$. We can find this by solving 

\begin{equation}
    \cos\theta = x, \quad \sin\theta = y
\end{equation}

For $z/|z| = x + iy$. Note, the argument of a product is the sum of the arguments of each product term. 

\begin{theorem}
    De Moisre's formula states that 
    \begin{equation}
        (\cos\theta + i\sin\theta)^n = \cos n\theta + i\sin n\theta
    \end{equation}
    Proof: use $e^{i\theta} = \cos\theta + i\sin\theta$. 
\end{theorem}

We can use De Moisre's formula to derive useful trigonometric identities, and matching the imaginary part with the other imaginary part, and etc. 

\section*{Roots of complex numbers}

Let $z = re^{i\theta}$, then we want to find 
\begin{equation}
    w^n = z
\end{equation}

Here, then 

\begin{align}
    w^n &= re^{i\theta}
\end{align}

Let $w = \rho e^{i\phi}$, then 

\begin{equation}
    \rho^n = r, \quad n\phi = \theta + 2\pi k,\; k \in \mathbb{Z}
\end{equation}

Clearly, $\rho = r^{1/n}$, and 

\begin{equation}
    \phi = \theta/n + 2\pi k/n
\end{equation}

Therefore, 

\begin{equation}
    z^{1/n} = \{r^{1/n} \exp(i(\frac{\theta}{n} + \frac{2\pi k}{n})): \quad k = 0, \dots, n-1\}
\end{equation}

There are many values of this number.

\begin{definition}
    Let $z \in \mathbb{C}$, then let $\theta = \arg(z)$ (principal value of $\arg(z)$), then 
    \begin{equation}
        r^{1/n}e^{i\arg(z)/n}
    \end{equation}
    This is the \underline{principal root}, with the notation of $^n\sqrt{z}$ (as opposed to $z^{1/n}$). 
\end{definition}

Suppose $z \neq 0$ and $z = re^{i\theta}$, then 

\begin{align}
    z^{1/2} = \{c_0, c_1\} \iff \{\sqrt{r}e^{i\theta/2}, \sqrt{r}e^{i(\theta/2 + 2\pi/2) \iff \sqrt{r}}e^{i\theta/2}e^{i\pi} = -\sqrt{r}e^{i\theta/2}\}
\end{align}

As another example, 

\begin{equation}
    1^{1/2} = \{1, -1\}
\end{equation}

However, $\sqrt{1}$ (principal root) is $1$. 

\subsection*{Example}

Evaluate $(-16)^{1/4}$ and find the principal value. Here, $\pi = \arg(-16)$. Then, 

\begin{equation}
    c_k = 2e^{(\pi/4 + 2\pi k/4 )i}
\end{equation}

Calcuate 

\begin{align}
    c_0 &= 2\left( \frac{\sqrt{2}}{2} + i\frac{\sqrt{2}}{2} \right) \\
    c_1 &= 2\left( -\frac{\sqrt{2}}{2} + i\frac{\sqrt{2}}{2} \right) \\
    c_2 &= 2\left( -\frac{\sqrt{2}}{2} - i\frac{\sqrt{2}}{2} \right) \\
    c_3 &= 2\left( \frac{\sqrt{2}}{2} - i\frac{\sqrt{2}}{2} \right)
\end{align}

Here, then the principal root is $c_0 = \sqrt{2} + i\sqrt{2}$ because $\pi = \arg(z)$. \textbf{This will probably be on the midterm.}



\end{document}